\section{Relação Com o Vídeo}

O conteúdo audiovisual produzido para este trabalho, com duração de 20 minutos, foi concebido como um recurso pedagógico complementar que visa traduzir a fundamentação algébrica deste documento em representações geométricas e dinâmicas. A proposta de complementaridade estabelece que enquanto este documento foca no rigor da modelagem matemática e na dedução dos resultados, o vídeo prioriza a intuição espacial e a visualização dos fenômenos descritos.

A relação entre os dois formatos é estruturada conforme os objetivos de aprendizagem de cada etapa:

\begin{itemize}
    \item \textbf{Contextualização e Diagnóstico Térmico:} A transição entre os desafios da computação clássica e a necessidade da integral dupla como "raio-X" térmico é apresentada no vídeo através de mapas de calor dinâmicos. Enquanto o texto detalha a fórmula da temperatura média, o vídeo permite visualizar o comportamento dos \textit{hotspots} sobre a área do processador.
    
    \item \textbf{Visualização da Transição de Escalas:} A mudança de paradigma do chip para o Qubit, discutida teoricamente na transição entre as cenas, é reforçada por recursos visuais de zoom e alterações de trilha sonora, que auxiliam o espectador a compreender a mudança de escala do domínio clássico para o quântico.
    
    \item \textbf{Modelagem de Incerteza Quântica:} A integral tripla e o cálculo do volume de erro são apresentados no vídeo mediante a visualização 3D da Esfera de Bloch e da nuvem de incerteza. Esta representação torna intuitiva a dependência cúbica entre a temperatura e a probabilidade de falha, consolidando os resultados numéricos obtidos neste documento.
    
    \item \textbf{Síntese da Relação Térmica-Quântica:} A correlação final entre a temperatura média e a probabilidade de erro ($P_{erro} \times T_{\text{média}}$) é consolidada no vídeo através de gráficos comparativos que demonstram visualmente como pequenas reduções térmicas impactam a fidelidade quântica, servindo como suporte visual à tabela de resultados aqui exposta.
\end{itemize}

Dessa forma, o vídeo funciona como uma interface de visualização para os modelos descritos, possuindo autonomia própria como material de divulgação científica, ao mesmo tempo em que aprofunda a compreensão dos cálculos apresentados neste trabalho escrito.